\documentclass[11pt]{article}

\usepackage{mathtools,amssymb,amsthm}
\usepackage[margin=1in]{geometry}
\usepackage{microtype}
\usepackage{url}

\raggedbottom

\setlength{\abovedisplayskip}{8pt}
\setlength{\belowdisplayskip}{8pt}
\setlength{\abovedisplayshortskip}{6pt}
\setlength{\belowdisplayshortskip}{6pt}

\linespread{0.97}

\title{A Finite 3-Adic Obstruction to Odd Collatz Cycles}
\author{Juan R. Marchionatto}
\date{}

% Theorem environments
\theoremstyle{definition}
\newtheorem{definition}{Definition}

\theoremstyle{plain}
\newtheorem{theorem}{Theorem}[section]
\newtheorem{lemma}[theorem]{Lemma}
\newtheorem{proposition}[theorem]{Proposition}

\theoremstyle{remark}
\newtheorem{remark}[theorem]{Remark}

\begin{document}
\maketitle

\begin{abstract}
We study the odd-to-odd dynamics of the Collatz map via the associated
exponent sequence defined by $3n+1=2^{e(n)}m$.
Finite windows of consecutive odd steps satisfy exact algebraic identities
which induce compatibility constraints modulo powers of $3$.
Under the hypothesis of an odd Collatz cycle, these constraints must hold
simultaneously for all circular windows of fixed length.

We encode these constraints as finite directed graphs.
A complete enumeration of directed cycles in the graphs corresponding to
window lengths $k=3$ and $k=4$ shows that no nonconstant exponent configuration
satisfies the required $3$-adic compatibility conditions.
The remaining constant exponent case is then eliminated by a direct algebraic
argument.

This yields a finite obstruction to the existence of nontrivial odd
Collatz cycles.
\end{abstract}

%--------------------------------
\section{Introduction}

The Collatz map sends a positive integer $n$ to $n/2$ if $n$ is even and
to $3n+1$ if $n$ is odd.
Despite its elementary formulation, the global behavior of this iteration
remains unresolved. The problem has been extensively studied in the literature;
see for example \cite{Lagarias1985,Lagarias2010}.
One of the fundamental open questions is whether nontrivial cycles exist.

It is natural to restrict attention to the induced odd-to-odd dynamics.
Each odd step can be written uniquely in the form
\[
3n+1 = 2^{e(n)} m,
\]
where $m$ is odd and $e(n)=v_2(3n+1)$.
Iterating this representation encodes odd Collatz trajectories as sequences
of exponents.

The central observation of this work is that finite windows of consecutive
odd steps satisfy exact algebraic identities.
When reduced modulo powers of $3$, these identities impose compatibility
constraints on ordered blocks of exponents that must hold simultaneously
under all circular shifts along a hypothetical cycle.

These circular compatibility conditions are encoded as finite directed graphs.
At window length three, the resulting $3$-adic constraints force the exponent
sequence to be constant modulo $18$.
At window length four, a refinement modulo $81$ eliminates all nonconstant
configurations.
The remaining constant case is resolved by a direct algebraic argument.

The argument is entirely finite and local in nature. It relies only on exact
window identities and exhaustive compatibility checks on finite graphs,
without invoking global convergence assumptions.

%--------------------------------
\section{Statement of the Main Result}\label{sec:main}

Let $T\colon\mathbb{N}\to\mathbb{N}$ denote the Collatz map,
\[
T(n)=
\begin{cases}
n/2 & \text{if } n \text{ is even},\\
3n+1 & \text{if } n \text{ is odd}.
\end{cases}
\]

We restrict attention to the induced odd-to-odd dynamics.
For an odd integer $n$, define
\[
3n+1 = 2^{e(n)} m,
\]
where $m$ is odd and $e(n)=v_2(3n+1)$.

An \emph{odd Collatz cycle} is a finite sequence of odd integers
\[
n_0,n_1,\dots,n_{p-1}
\]
such that
\[
3n_i+1=2^{e_i}n_{i+1}
\quad (i=0,\dots,p-1),
\]
with indices taken modulo $p$.

The trivial cycle consists of the single fixed point $n_0=1$.

\begin{theorem}
There exist no nontrivial odd Collatz cycles.
\end{theorem}

%--------------------------------
\section{Algebraic Window Identity}\label{sec:algebraic-w-ident}

Let $(n_i)$ be an odd Collatz trajectory and $(e_i)$ its exponent sequence,
\[
3n_i+1 = 2^{e_i} n_{i+1}.
\]

\subsection{Iterated identity}

Fix an integer $k\ge 1$.
By repeated substitution, one obtains the exact identity
\begin{equation}\label{eq:window_identity}
3^k n_{i+k}
=
2^{E_i} n_i
-
\sum_{j=0}^{k-1}
3^{j} \,
2^{\,e_{i+j+1}+\cdots+e_{i+k-1}},
\end{equation}
where
\[
E_i = e_i+e_{i+1}+\cdots+e_{i+k-1}.
\]

\subsection{Reduction modulo $3^k$}

Reducing \eqref{eq:window_identity} modulo $3^k$ eliminates the left-hand side and yields
\begin{equation}\label{eq:window_congruence}
2^{E_i} n_i
\equiv
\sum_{j=0}^{k-1}
3^{j} \,
2^{\,e_{i+j+1}+\cdots+e_{i+k-1}}
\pmod{3^k}.
\end{equation}

Since $\gcd(2,3^k)=1$, the factor $2^{E_i}$ is invertible modulo $3^k$.
Thus we may rewrite \eqref{eq:window_congruence} in the form
\begin{equation}\label{eq:Rk_definition}
n_i \equiv R_k(e_i,\dots,e_{i+k-1})
\pmod{3^k},
\end{equation}
where
\[
R_k(e_i,\dots,e_{i+k-1})
=
2^{-E_i}
\sum_{j=0}^{k-1}
3^{j} \,
2^{\,e_{i+j+1}+\cdots+e_{i+k-1}}
\pmod{3^k}.
\]

\subsection{Consequence for cycles}

If $(n_0,\dots,n_{p-1})$ is an odd Collatz cycle, then for every
$i\in\mathbb{Z}/p\mathbb{Z}$ the congruence \eqref{eq:Rk_definition}
holds for the window
\[
(e_i,e_{i+1},\dots,e_{i+k-1}),
\]
with indices taken modulo $p$.

%--------------------------------
\section{Circular Compatibility as a Finite Graph}\label{sec:circ-compat}

Fix $k \ge 1$.

For each ordered $k$-tuple of exponents
\[
(e_i,\dots,e_{i+k-1}),
\]
define the residue
\[
R_k(e_i,\dots,e_{i+k-1})
\in \mathbb{Z}/3^k\mathbb{Z}
\]
as in \eqref{eq:Rk_definition}.

\subsection{Compatibility condition}

Along an odd trajectory we have
\[
3n_i + 1 = 2^{e_i} n_{i+1}.
\]

Reducing modulo $3^k$ and substituting the window residues yields
\begin{equation}\label{eq:compatibility}
R_k(e_{i+1},\dots,e_{i+k})
\equiv
\bigl(3R_k(e_i,\dots,e_{i+k-1}) + 1\bigr)\,
2^{-e_i}
\pmod{3^k}.
\end{equation}

\subsection{Directed graph encoding}

For each fixed window length $k$, we associate a directed graph $G_k$ as follows:

\begin{itemize}
\item Vertices: all ordered $k$-tuples of exponent residues modulo
the multiplicative order of $2$ modulo $3^k$.
\item Directed edge:
\[
(e_0,\dots,e_{k-1}) \longrightarrow (e_1,\dots,e_k)
\]
whenever the compatibility condition \eqref{eq:compatibility} holds.
\end{itemize}

\begin{lemma}\label{lem:cycle_to_Gk}
For each fixed $k$, if an odd Collatz cycle exists, then its length-$k$
exponent windows form a directed cycle in the associated graph $G_k$.
\end{lemma}

\begin{proof}
In a cycle of length $p$, every length-$k$ window occurs,
including those wrapping around the cycle.
Each must satisfy \eqref{eq:compatibility}.
Therefore the sequence of windows produces a directed cycle in $G_k$.
\end{proof}

%--------------------------------
\section{Finite Obstruction at Window Length $k=3$}\label{sec:k3}

We specialize the window identity to length $k=3$ and work modulo $3^3=27$.
Since $2$ is a unit modulo $27$ and powers of $2$ are periodic modulo $27$,
only exponent residues modulo the multiplicative order of $2$ mod $27$
are relevant. For $k\ge 1$, the multiplicative order of $2$ modulo $3^k$ is
\[
\operatorname{ord}_{3^k}(2)=2\cdot 3^{k-1},
\qquad\text{hence}\qquad
\operatorname{ord}_{27}(2)=18.
\]
Therefore, every residue condition modulo $27$ depends only on exponents
reduced modulo $18$.

\begin{lemma}\label{lem:mod18_reduction}
For $k=3$, the residue map modulo $27$ depends only on exponent classes modulo $18$.
More precisely, the quantity
\[
R_3(a,b,c)
=
2^{-(a+b+c)}
\Bigl(2^{b+c}+3\cdot 2^{c}+9\Bigr)
\pmod{27}
\]
is determined entirely by the classes of $a,b,c$ in $\mathbb{Z}/18\mathbb{Z}$.
Moreover, the compatibility condition
\[
R_3(b,c,d)\equiv \bigl(3R_3(a,b,c)+1\bigr)\,2^{-a}\pmod{27}
\]
also depends only on the classes of $a,b,c,d$ modulo $18$.
\end{lemma}

\begin{proof}
Since $\operatorname{ord}_{27}(2)=18$, we have
\[
2^{x+18m}\equiv 2^x \pmod{27}
\qquad\text{for all }x,m\in\mathbb{Z}.
\]
Because $2$ is invertible modulo $27$, the same holds for inverses:
\[
2^{-(x+18m)}\equiv 2^{-x}\pmod{27}.
\]
Thus every power of $2$ and every inverse power of $2$ appearing in the definition
of $R_3(a,b,c)$ depends only on the corresponding exponent modulo $18$.
Therefore $R_3(a,b,c)$ itself depends only on $(a,b,c)\bmod 18$.
The same argument applies to the compatibility condition, since both
$R_3(a,b,c)$ and $2^{-a}$ depend only on the corresponding residue classes
modulo $18$.
\end{proof}

\subsection{The length--three residue map}

For a triple $(a,b,c)\in(\mathbb{Z}/18\mathbb{Z})^3$, define
\begin{equation}\label{eq:R3def}
R_3(a,b,c)
:=
2^{-(a+b+c)}
\Bigl(2^{b+c}+3\cdot 2^{c}+9\Bigr)
\;\in\;\mathbb{Z}/27\mathbb{Z},
\end{equation}
where all powers of $2$ and inverses are taken modulo $27$.
This is well-defined since $2$ is invertible modulo $27$.

If $(e_i,e_{i+1},e_{i+2})$ is a realizable length--three window along an odd
Collatz trajectory and we reduce exponents modulo $18$, then the window identity
implies
\[
n_i \equiv R_3(e_i,e_{i+1},e_{i+2}) \pmod{27}.
\]

\subsection{Circular compatibility as a directed graph}

Along an odd trajectory we have
\[
3n_i+1 = 2^{e_i}n_{i+1}.
\]
Reducing modulo $27$ and substituting the window residues gives the compatibility
constraint
\begin{equation}\label{eq:compat3}
R_3(e_{i+1},e_{i+2},e_{i+3})
\equiv
\bigl(3\,R_3(e_i,e_{i+1},e_{i+2})+1\bigr)\,2^{-e_i}
\pmod{27},
\end{equation}
with all $e_j$ understood modulo $18$.

Define a directed graph $G_3$ whose vertex set is $(\mathbb{Z}/18\mathbb{Z})^3$.
We draw a directed edge
\[
(a,b,c)\longrightarrow (b,c,d)
\]
if and only if the compatibility congruence \eqref{eq:compat3} holds.

\begin{lemma}[Real cycles project to $G_3$]\label{lem:projection_to_G3}
Let
\[
(n_0,n_1,\dots,n_{p-1})
\]
be an odd Collatz cycle with exponent sequence $(e_i)$.
Then the sequence of triples
\[
(e_i,e_{i+1},e_{i+2}) \bmod 18
\]
forms a directed cycle in $G_3$.
\end{lemma}

\begin{proof}
For every index $i$ along the cycle, the length--three window identity
implies
\[
n_i \equiv R_3(e_i,e_{i+1},e_{i+2}) \pmod{27}.
\]

Since the window identity is an exact algebraic identity for every odd
Collatz trajectory, this congruence holds for every index $i$ along the
cycle.

The odd-to-odd recurrence gives
\[
3n_i+1 = 2^{e_i} n_{i+1}.
\]

Reducing modulo $27$ yields
\[
n_{i+1} \equiv (3n_i+1)\,2^{-e_i}\pmod{27}.
\]

Substituting the window residues gives
\[
R_3(e_{i+1},e_{i+2},e_{i+3})
\equiv
\bigl(3R_3(e_i,e_{i+1},e_{i+2})+1\bigr)\,2^{-e_i}
\pmod{27}.
\]

By Lemma~\ref{lem:mod18_reduction}, this condition depends only on the
residue classes modulo $18$.
Therefore the triples
\[
(e_i,e_{i+1},e_{i+2}) \bmod 18
\]
follow the edge rule defining $G_3$.

Since the cycle is circular, these triples form a directed cycle in $G_3$.
\end{proof}

\subsection{Consequence: rigidity modulo $18$}

The key finite obstruction at window length three is the following.

\begin{lemma}\label{lem:G3_only_constant_cycles}
The directed graph $G_3$ has no directed cycles other than the constant
triples modulo $18$
\[
(a,a,a)\to(a,a,a),
\qquad a\in\mathbb{Z}/18\mathbb{Z}.
\]
\end{lemma}

Recall that the vertices of $G_3$ are triples of exponent residues
\[
(a,b,c)\in(\mathbb{Z}/18\mathbb{Z})^3.
\]
Thus a vertex $(a,a,a)$ represents three consecutive exponents that are
equal modulo $18$, not necessarily equal as integers.

Lemma~\ref{lem:G3_only_constant_cycles} is a finite statement about a
directed graph with $18^3$ vertices. Once it holds, any exponent sequence
arising from a hypothetical odd Collatz cycle must satisfy
\[
e_i \equiv e_{i+1} \equiv \cdots \equiv e_{i+p-1} \pmod{18}.
\]
Indeed, if a cycle existed, the sequence of triples
\[
(e_i,e_{i+1},e_{i+2}) \bmod 18
\]
would form a directed cycle in $G_3$. By the lemma, the only possible
directed cycles are the constant triples $(a,a,a)$, which implies that
all exponents along the cycle lie in the same residue class modulo $18$.

In particular, the exponent sequence of any hypothetical odd Collatz cycle
lies in a single residue class modulo $18$.

\medskip
A computational verification of Lemma~\ref{lem:G3_only_constant_cycles} is
given in \ref{app:computation} by explicit construction of $G_3$ and enumeration of
its strongly connected components.

%--------------------------------
\section{Refinement at Window Length $k=4$}\label{sec:k4}

By Lemma~\ref{lem:G3_only_constant_cycles}, any exponent sequence arising
from a hypothetical odd Collatz cycle satisfies
\[
e_i \equiv r \pmod{18}
\]
for some fixed residue $r \in \mathbb{Z}/18\mathbb{Z}$.

We therefore write
\[
e_i = r + 18 t_i,
\qquad t_i \in \mathbb{Z}.
\]

\subsection{Reduction modulo $81$}

Specializing the window identity to length $k=4$ and reducing modulo
$3^4 = 81$, we analyze the resulting compatibility condition.

Since
\[
e_i = r+18t_i,
\]
we have
\[
2^{e_i}=2^r(2^{18})^{t_i}\pmod{81}.
\]
Now
\[
2^{18}\equiv 28\pmod{81},
\qquad
28^3\equiv 1\pmod{81},
\]
so the residue of $(2^{18})^{t_i}$ modulo $81$ depends only on $t_i\bmod 3$.
Moreover, every power of $2$ appearing in the length--four window identity
is obtained from sums of exponents of the form
\[
e_{i_1}+\cdots+e_{i_m}
= mr + 18(t_{i_1}+\cdots+t_{i_m}),
\]
and therefore also depends only on the corresponding residue classes
$t_{i_j}\bmod 3$.
Hence all powers of $2$ occurring in the length--four residue map, and
therefore the entire compatibility condition modulo $81$, depend only on
\[
t_i \in \mathbb{Z}/3\mathbb{Z}.
\]

Define the residue map
\begin{equation}\label{eq:R4def}
R_4(a,b,c,d)
=
2^{-(\epsilon(a)+\epsilon(b)+\epsilon(c)+\epsilon(d))}
\Bigl(
2^{\epsilon(d)+\epsilon(c)+\epsilon(b)}
+ 3\,2^{\epsilon(d)+\epsilon(c)}
+ 9\,2^{\epsilon(d)}
+ 27
\Bigr)
\pmod{81},
\end{equation}
where
\[
\epsilon(x)=r+18x,
\qquad x\in\mathbb{Z}/3\mathbb{Z}.
\]

\subsection{Directed graph $G_r$}

Define a directed graph $G_r$ as follows:

\begin{itemize}
\item Vertices: all quadruples $(a,b,c,d)\in(\mathbb{Z}/3\mathbb{Z})^4$.
\item Directed edge:
\[
(a,b,c,d)\longrightarrow (b,c,d,x)
\]
if and only if the compatibility condition
\[
R_4(b,c,d,x)
\equiv
\bigl(3R_4(a,b,c,d)+1\bigr)\,2^{-\epsilon(a)}
\pmod{81}
\]
holds.
\end{itemize}

\begin{lemma}\label{lem:cycle_implies_Gr_cycle}
If an odd Collatz cycle exists with exponent residue $r \pmod{18}$,
then the sequence of quadruples
\[
(t_i,t_{i+1},t_{i+2},t_{i+3})
\]
forms a directed cycle in $G_r$.
\end{lemma}

\begin{proof}
As in the length--three case, every length--four window along the cycle,
including those wrapping around the closing identification, must satisfy
the window identity modulo $81$.
Since the window identity is an exact algebraic identity for every odd
Collatz trajectory, the corresponding residue relation holds at every
index of the cycle.
Therefore consecutive windows satisfy the compatibility condition defining
$G_r$, and hence produce a directed cycle.
\end{proof}
\subsection{Elimination of nonconstant configurations}
The key finite obstruction at window length four is the following.

\begin{lemma}\label{lem:Gr_only_constant_cycles}
For each $r\in\mathbb{Z}/18\mathbb{Z}$, the directed graph $G_r$
has no directed cycles other than the constant ones
\[
(a,a,a,a)\to(a,a,a,a),
\qquad a\in\mathbb{Z}/3\mathbb{Z}.
\]
\end{lemma}

Lemma~\ref{lem:Gr_only_constant_cycles} is a finite statement about a directed graph
with $3^4=81$ vertices.
A computational verification is provided in \ref{app:computation}.

Consequently,
\[
t_0=t_1=\cdots=t_{p-1},
\]
and therefore
\[
e_0=e_1=\cdots=e_{p-1}.
\]

Hence any hypothetical odd Collatz cycle must have constant exponent sequence.

%--------------------------------
\section{The Constant Exponent Case}\label{sec:const}

By Section~\ref{sec:k4}, any hypothetical odd Collatz cycle must have constant exponent
sequence. Let
\[
e_i = e
\quad\text{for all } i.
\]

\subsection{Global cycle identity}

Composing $p$ consecutive odd steps yields the identity
\[
n_{i+p}
=
\frac{2^{pe}}{3^p}\, n_i
-
\frac{1}{3^p}
\sum_{j=0}^{p-1}
3^{j}\, 2^{(p-1-j)e}.
\]

If the trajectory forms a cycle of length $p$, then $n_{i+p}=n_i$.
Rearranging gives
\begin{equation}\label{eq:constant_cycle_identity}
n_i\bigl(2^{pe}-3^p\bigr)
=
\sum_{j=0}^{p-1}
3^{j}\, 2^{(p-1-j)e}.
\end{equation}

Solving for $n_i$ yields
\begin{equation}\label{eq:constant_solution}
n_i
=
\frac{\displaystyle\sum_{j=0}^{p-1}3^{j}2^{(p-1-j)e}}
{2^{pe}-3^p}.
\end{equation}

\subsection{Case analysis}

\paragraph{Case $e=1$.}

Then $2^{p}-3^p<0$ for all $p\ge1$, while the numerator in
\eqref{eq:constant_solution} is positive.
Thus $n_i<0$, which is impossible.

\paragraph{Case $e\ge3$.}

Define the odd-to-odd map
\[
T_e(n)=\frac{3n+1}{2^e}.
\]

For $e\ge3$ we have $2^e-3\ge5$.
For every $n\ge1$,
\[
3n+1 < 2^e n
\quad\Longleftrightarrow\quad
1 < (2^e-3)n.
\]
Since $(2^e-3)n\ge5n\ge5$, this inequality holds.
Hence
\[
T_e(n)<n.
\]

Thus the odd terms strictly decrease at each step,
which is incompatible with the existence of a cycle.

\paragraph{Case $e=2$.}

In this case
\[
T_2(n)=\frac{3n+1}{4}.
\]
Equation \eqref{eq:constant_solution} yields $n_i=1$ for all $p\ge1$.
This corresponds to the trivial odd cycle.

\subsection{Conclusion}

The analysis above shows that if the exponent sequence of an odd Collatz
cycle is constant, then the only possible cycle is the trivial one
corresponding to $n=1$.

%--------------------------------
\begin{proof}[Proof of the Main Theorem]
Assume a nontrivial odd Collatz cycle exists.

By Section~\ref{sec:k3}, the exponent sequence must be constant modulo $18$.
By Section~\ref{sec:k4}, it must therefore be constant.
Section~\ref{sec:const} shows that the constant exponent case produces
only the trivial cycle.

This contradicts the assumption that a nontrivial odd Collatz cycle exists.
\end{proof}

%--------------------------------
\section*{Conclusion}

We have shown that the existence of an odd Collatz cycle would force the
associated exponent sequence to satisfy strong local compatibility
conditions arising from the window identities.
Encoding these conditions as finite directed graphs yields a complete
finite obstruction, forcing the exponent sequence to be constant.
The constant case produces only the trivial cycle.

Hence no nontrivial odd Collatz cycles exist.

%--------------------------------
\appendix
\renewcommand{\thesection}{Appendix \Alph{section}}
\renewcommand{\thesubsection}{\Alph{section}.\arabic{subsection}}

\section{Computational Verification of the Graph Constraints}\label{app:computation}

The statements of Lemma~\ref{lem:G3_only_constant_cycles}
and Lemma~\ref{lem:Gr_only_constant_cycles} reduce to finite
computations on directed graphs.
\subsection{Construction of $G_3$}
The graph $G_3$ has vertex set
\[
(\mathbb{Z}/18\mathbb{Z})^3,
\]
so it contains $18^3=5832$ vertices.

For each vertex $(a,b,c)$ and each $d\in\mathbb{Z}/18\mathbb{Z}$,
we test whether the compatibility condition
\[
R_3(b,c,d)
\equiv
(3R_3(a,b,c)+1)\,2^{-a}
\pmod{27}
\]
holds. If so, we add the directed edge
\[
(a,b,c)\rightarrow(b,c,d).
\]

The resulting directed graph has $5832$ vertices and $5832$ edges.
In particular, each vertex has out-degree exactly one.

\subsection{Cycle detection}

To detect cycles, we compute the strongly connected components (SCCs)
of the directed graph using Tarjan's algorithm.
The computation shows that the only SCCs containing directed cycles
are the trivial self-loops
\[
(a,a,a)\to(a,a,a)
\]
for $a\in\mathbb{Z}/18\mathbb{Z}$.

In particular, $G_3$ contains no nonconstant directed cycles.

\subsection{Verification for $k=4$}

For the refinement step, the graph $G_r$ has vertex set
\[
(\mathbb{Z}/3\mathbb{Z})^4,
\]
containing $3^4=81$ vertices.
This verification is carried out separately for each residue
\[
r\in\mathbb{Z}/18\mathbb{Z},
\]
so in total $18$ directed graphs are checked.
Edges are defined by the compatibility condition modulo $81$.

The same SCC analysis shows that the only directed cycles
are the constant ones
\[
(a,a,a,a)\to(a,a,a,a),
\qquad a\in\mathbb{Z}/3\mathbb{Z}.
\]

\subsection{Implementation}

The computations were implemented in Python. The source code used to
construct the graphs and enumerate their strongly connected components is
publicly available at

\url{https://github.com/jmarchionatto/collatz-3adic-cycle-obstruction}.

%--------------------------------
\section*{AI Assistance Disclosure}

This work was developed with assistance from large language models for drafting,
structuring, and refining exposition, as well as for checking algebraic derivations.
The mathematical ideas, results, and overall research direction are due to the
author, who takes full responsibility for the content.

\begin{thebibliography}{9}

\bibitem{Lagarias1985}
J.~C.~Lagarias,
\newblock The $3x+1$ problem and its generalizations,
\newblock \emph{Amer. Math. Monthly} \textbf{92} (1985), no.~1, 3--23.

\bibitem{Lagarias2010}
J.~C.~Lagarias,
\newblock The $3x+1$ problem: An annotated bibliography (1963--1999),
\newblock \emph{arXiv preprint} arXiv:math/0309224, revised 2010.

\end{thebibliography}
\end{document}